\section{Benchmark Motivation and Construction}
\label{sec:benchmark}

\subsection{Scope and Definitions}
% Original scaffold TODO: Define the cultural knowledge and video understanding
% abilities evaluated by the benchmark, including what is outside its scope.
We study how models interpret cultural elements when their surrounding
context does not point to a single cultural association. A performance can
retain a recognizable form while its performers, music, or social purpose
change. The same problem arises when a garment combines a familiar shape
with another textile tradition, or when a food practice enters a different
ceremonial setting. Recognizing each element and explaining its role are
separate requirements.

\paragraph{Culture and cultural association.}
We use \emph{culture} operationally to refer to learned and shared practices,
material expressions, and meanings associated with communities and
traditions. This definition supports annotation of observable examples; it
does not treat a culture as a fixed list of national traits. A
\emph{cultural association} links an observed element to a documented
tradition, community, region, or cultural context. Several associations may
be valid. A national label does not establish exclusive ownership,
historical origin, or the identity of the people in a video.
A carrier or receiving context may have a documented cultural association,
but appearance alone is not evidence of a person's affiliation.
\todo{Confirm the cultural descriptors and references used to validate associations.}

\paragraph{Cue, component, and evidence.}
A \emph{cultural cue} is a particular element in a video that supports a
cultural interpretation. Its \emph{component type} describes what the
element is, such as a practice, movement, musical element, garment,
language, setting, or social interaction. Its \emph{evidence channel}
records where the supporting information appears: visual appearance,
visual action, non-speech audio, speech, or visible text. For example, a
musical practice is a component, while the sound recording is an evidence
channel. Expert knowledge may be needed to connect the observed element
to a tradition; such knowledge should be documented separately from the
observation itself.

\paragraph{Cultural mixing.}
We define a cultural-mixing example as a scene or coherent event containing
distinguishable elements associated with different cultural contexts.
These may be two cultural cues, or a cue and a separately documented
carrier or receiving context. A change in performers or setting can
therefore qualify without changing the form itself. This includes
coexistence, component incorporation, and adaptation or fusion.
Technical fusion is therefore not required. However, a compilation that
places unrelated clips next to each other does not by itself establish a
mixed event. Nor does assigning several valid cultural labels to one cue
automatically create multiple cue instances.
The sample's inclusion and a task's available evidence are separate:
claims requiring external provenance are evaluated only when that
information is supplied in a declared metadata-assisted condition.
\todo{Validate these inclusion boundaries on examples with shared traditions and overlapping contexts.}

\paragraph{What the benchmark measures.}
In this paper, \emph{cultural understanding} means recognizing culturally
associated elements, identifying evidence for them, and, in the candidate
relation task, assigning their roles within the observed event. We do not
measure a model's cultural identity, values, authenticity judgments, or
ability to determine cultural ownership. Historical influence and creator
intent are outside the primary V and VA tasks unless separate evidence is given.

% The supplied world-modelling guidelines motivate attention to activities,
% local practices, and overlapping region/language labels. Their 8--15 second
% limit, prefix/suffix prediction, and single-category rule are not adopted.
The working scope extends beyond the initial performing-arts search pilot
to food practices, clothing and materials, crafts, architecture and
decoration, rituals, and media adaptations.
\todo{Confirm the final domain list, cultural coverage, and feasible collection targets.}

\subsection{Data Collection}
% Original scaffold TODO: Describe video sources, inclusion criteria,
% cultural coverage, licensing, and filtering.
The following pipeline is a proposed construction protocol. The retrieved
pilot links establish candidate sources, not a completed benchmark or a
verified estimate of collection yield.

\paragraph{Sources and retrieval.}
The pilot uses publicly discoverable YouTube videos. We propose searching
for combinations of a cultural form or component and a second component
or context. Queries may use local names, transliterations, and the
language of the receiving community. For example, a practice name can be
combined with a local event or institution. Model assistance may suggest
queries and related entities, but a retrieved title or model judgment
does not establish a gold label. Each candidate should retain its source
URL, video identifier, retrieval date, exact query, query language, and
any parent query used for expansion.
\todo{Confirm source platforms, retrieval budgets, search models, and the saved query log.}

\paragraph{Screening and evaluation units.}
Reviewers should inspect the audiovisual content to confirm that the
proposed cues are observable and belong to a coherent scene or event. They should
also record whether the cultural interpretation depends on a title,
description, subtitle, or other external information. A video's uploader
location alone is insufficient evidence of a cultural context. Source
videos and evaluation clips are different units: a long source may
contain several relevant intervals. We propose retaining the source
identifier and original timestamps when extracting an event, and storing
discontinuous cue evidence as multiple intervals.
\todo{Set source-duration filters, evaluation-clip limits, resolution requirements, and repeated-event handling.}

\paragraph{Rights and duplication.}
Public accessibility does not establish permission to redistribute a
video. The source record should include the license or permission status
and the intended release format. Duplicate video identifiers should be
merged. Reuploads, overlapping clips, and recordings of the same event
should be grouped before any data split, with manual review of suspected
matches.
\todo{Confirm the rights-review process, redistribution policy, duplicate-detection method, and split grouping rules.}

\subsection{Annotation and Quality Control}
% Original scaffold TODO: Describe annotator backgrounds, annotation stages,
% agreement checks, and adjudication.
We propose separating a factual scene description from a cultural note.
The first records visible or audible events without assigning unobserved
intent. The second explains the cultural association and records its
supporting source or relevant expertise. This separation makes it
possible to audit a correct cultural claim that lacks evidence in the
particular video.

For each retained event, annotators would identify cue instances, assign
component types and cultural associations, and mark supporting evidence
channels and intervals. They would then label candidate relations between
specified cues and their participating entities or contexts. Relation annotations may contain more than one edge
per event; a single forced video-level category would hide valid
combinations. The task definitions and proposed role inventory appear in
Section~\ref{sec:tasks}.

Independent annotation should precede adjudication. Annotators should
have relevant knowledge of the traditions they label, with a way to flag
uncertainty and request additional review. The original search query and
retrieval label should be hidden during independent annotation to reduce
confirmation bias. Unresolved associations or roles should be retained
as uncertain rather than silently converted into a single gold answer.
\todo{Confirm recruitment, cultural expertise, annotation languages, compensation, number of judgments per item, and adjudication procedure.}

Quality control should report agreement separately for cue recognition,
cultural association, evidence channels, relation frames, and role
assignment. A pilot should test both clear examples and boundary cases,
including whether embedding and realization can be distinguished
reliably. A relation class should enter the main evaluation only after
its instructions produce reproducible judgments.
\todo{Choose agreement measures and acceptance criteria; report measured agreement after the pilot.}

\subsection{Dataset Statistics}
\todo{Report verified dataset statistics and distributions.}
The final report should distinguish retrieved sources, screened sources,
accepted events, released clips, cue instances, and task items. It should
also report duration, domain and cultural coverage, evidence-channel
combinations, relation frequencies, unresolved cases, rejection reasons,
and rights status. Counts from retrieval should not be presented as
counts of annotated examples.
